% GENERATED FILE. Edit canonical lessons, then run npm run generate:books.
% canonical-source-hash: fnv1a64:6c14ef4d5dec4f51
% canonical-lessons: ES-C621-impedir, ES-C621-facilitar, ES-C621-dificultar, ES-C621-exento, ES-C621-prohibido

\chapter{Lo hace posible o no}
\label{ch:es-lo-hace-posible-o-no}
\hlchaptermodality{621}

\begin{chapteropening}
\textbf{By the end of this chapter:} \emph{I can say what makes something easier or harder, what is forbidden and who is exempt.}

Complete the last lesson of chapter 621: I can say what makes something easier or harder, what is forbidden and who is exempt.
\end{chapteropening}

\section[impedir]{impedir — to stop, to prevent}

\label{lesson:ES-C621-impedir}



\begin{quote}
Before the new one: say the Spanish for complicated, then the Spanish for to raise.
\end{quote}

\subsection*{You'll want to know: impedir}
\textbf{impedir} — \textquotedblleft{}to stop, to prevent\textquotedblright{}.

\begin{quote}
La lluvia nos impidió salir.
\end{quote}

\begin{quote}
\emph{The rain stopped us going out.}
\end{quote}

Like \emph{pedir}: \emph{impido, impidió}.

\begin{cousinweb}
\emph{impedir} is a learned borrowing from Latin \emph{impedire} 'to hinder, to entangle the feet', built on \emph{pes, pedis} 'foot' (EtymDB).
\end{cousinweb}

\begin{grammarlens}[title={What you've built}]
One more word for telling, explaining and describing at B1.
\end{grammarlens}

\subsection*{Your turn}
\begin{itemize}
\raggedright
  \item \emph{Say it:} \emph{impedir}
  \item \emph{Say it:} \emph{impedir}, once more
  \item \emph{Say it:} a sentence of your own with \emph{impedir}
  \item \emph{Recall:} say the Spanish for complicated, then the Spanish for to raise, then say \emph{impedir} again
\end{itemize}

\begin{tcolorbox}[breakable,colback=teal!4,colframe=teal!35!black,title={Before you move on}]
What does impedir mean? (\textquotedblleft{}To stop, to prevent\textquotedblright{}.) Say it once more.
\end{tcolorbox}
\section[facilitar]{facilitar — to make easier; to provide}

\label{lesson:ES-C621-facilitar}



\begin{quote}
Before the new one: say the Spanish for to raise, then the Spanish for to stop.
\end{quote}

\subsection*{You'll want to know: facilitar}
\textbf{facilitar} — \textquotedblleft{}to make easier; to provide\textquotedblright{}.

\begin{quote}
El ordenador facilita el trabajo.
\end{quote}

\begin{quote}
\emph{The computer makes the work easier.}
\end{quote}

On forms: \emph{facilite sus datos} — provide your details.

\begin{cousinweb}
\emph{facilitar} is built on \emph{fácil}, a learned borrowing from Latin \emph{facilis} 'easy to do', from \emph{facere} 'to do' (EtymDB).
\end{cousinweb}

\begin{grammarlens}[title={What you've built}]
One more word for telling, explaining and describing at B1.
\end{grammarlens}

\subsection*{Your turn}
\begin{itemize}
\raggedright
  \item \emph{Say it:} \emph{facilitar}
  \item \emph{Say it:} \emph{facilitar}, once more
  \item \emph{Say it:} a sentence of your own with \emph{facilitar}
  \item \emph{Recall:} say the Spanish for to raise, then the Spanish for to stop, then say \emph{facilitar} again
\end{itemize}

\begin{tcolorbox}[breakable,colback=teal!4,colframe=teal!35!black,title={Before you move on}]
What does facilitar mean? (\textquotedblleft{}To make easier; to provide\textquotedblright{}.) Say it once more.
\end{tcolorbox}
\section[dificultar]{dificultar — to make difficult, to hinder}

\label{lesson:ES-C621-dificultar}



\begin{quote}
Before the new one: say the Spanish for to stop, then the Spanish for to make easier.
\end{quote}

\subsection*{You'll want to know: dificultar}
\textbf{dificultar} — \textquotedblleft{}to make difficult, to hinder\textquotedblright{}.

\begin{quote}
La nieve dificulta el viaje.
\end{quote}

\begin{quote}
\emph{The snow makes the journey difficult.}
\end{quote}

The opposite of \emph{facilitar}.

\begin{cousinweb}
\emph{dificultar} is a learned borrowing from Latin \emph{difficultare} 'to make hard', built on \emph{difficultas}, from \emph{facilis} 'easy' (EtymDB).
\end{cousinweb}

\begin{grammarlens}[title={What you've built}]
One more word for telling, explaining and describing at B1.
\end{grammarlens}

\subsection*{Your turn}
\begin{itemize}
\raggedright
  \item \emph{Say it:} \emph{dificultar}
  \item \emph{Say it:} \emph{dificultar}, once more
  \item \emph{Say it:} a sentence of your own with \emph{dificultar}
  \item \emph{Recall:} say the Spanish for to stop, then the Spanish for to make easier, then say \emph{dificultar} again
\end{itemize}

\begin{tcolorbox}[breakable,colback=teal!4,colframe=teal!35!black,title={Before you move on}]
What does dificultar mean? (\textquotedblleft{}To make difficult, to hinder\textquotedblright{}.) Say it once more.
\end{tcolorbox}
\section[exento]{exento — exempt, free from}

\label{lesson:ES-C621-exento}



\begin{quote}
Before the new one: say the Spanish for to make easier, then the Spanish for to make difficult.
\end{quote}

\subsection*{You'll want to know: exento}
\textbf{exento} — \textquotedblleft{}exempt, free from\textquotedblright{}.

\begin{quote}
Los niños están exentos de pagar.
\end{quote}

\begin{quote}
\emph{Children are exempt from paying.}
\end{quote}

\emph{Exento de} + noun or infinitive.

\begin{cousinweb}
\emph{exento} is a learned borrowing from Latin \emph{exemptus} 'taken out, freed', from \emph{eximere}, built on \emph{emere} 'to take' (Lewis \& Short).
\end{cousinweb}

\begin{grammarlens}[title={What you've built}]
One more word for telling, explaining and describing at B1.
\end{grammarlens}

\subsection*{Your turn}
\begin{itemize}
\raggedright
  \item \emph{Say it:} \emph{exento}
  \item \emph{Say it:} \emph{exento}, once more
  \item \emph{Say it:} a sentence of your own with \emph{exento}
  \item \emph{Recall:} say the Spanish for to make easier, then the Spanish for to make difficult, then say \emph{exento} again
\end{itemize}

\begin{tcolorbox}[breakable,colback=teal!4,colframe=teal!35!black,title={Before you move on}]
What does exento mean? (\textquotedblleft{}Exempt, free from\textquotedblright{}.) Say it once more.
\end{tcolorbox}
\section[prohibido]{prohibido — forbidden, prohibited}

\label{lesson:ES-C621-prohibido}



\begin{quote}
Before the new one: say the Spanish for to make difficult, then the Spanish for exempt.
\end{quote}

\subsection*{You'll want to know: prohibido}
\textbf{prohibido} — \textquotedblleft{}forbidden, prohibited\textquotedblright{}.

\begin{quote}
Prohibido aparcar.
\end{quote}

\begin{quote}
\emph{No parking.}
\end{quote}

Common on signs: \emph{prohibido el paso} — no entry.

\begin{cousinweb}
\emph{prohibido} is the participle of \emph{prohibir}, a learned borrowing from Latin \emph{prohibere} 'to hold back', built on \emph{habere} 'to hold' (EtymDB).
\end{cousinweb}

\begin{grammarlens}[title={What you've built}]
One more word for telling, explaining and describing at B1.
\end{grammarlens}

\subsection*{Your turn}
\begin{itemize}
\raggedright
  \item \emph{Say it:} \emph{prohibido}
  \item \emph{Say it:} \emph{prohibido}, once more
  \item \emph{Say it:} a sentence of your own with \emph{prohibido}
  \item \emph{Recall:} say the Spanish for to make difficult, then the Spanish for exempt, then say \emph{prohibido} again
\end{itemize}

\begin{tcolorbox}[breakable,colback=teal!4,colframe=teal!35!black,title={Before you move on}]
What does prohibido mean? (\textquotedblleft{}Forbidden, prohibited\textquotedblright{}.) Say it once more.
\end{tcolorbox}
